\section{Empirical Results}
\label{sec:results}

\subsection{Experimental Setup}

We compare Phase~1 \cvd{} and Phase~2 \cvdgeo{} on three states representing
the spectrum from regular geometry to severely coastal:

\begin{itemize}
\item \textbf{NC} (North Carolina, $k=14$, $n \approx 2{,}660$ tracts):
  regular inland geometry, competitive two-party state.
  BFS hop-count and Euclidean distance agree well; Phase~2 improvement expected
  to be modest.
\item \textbf{FL} (Florida, $k=28$, $n \approx 4{,}200$ tracts):
  the panhandle extends 200 miles westward; Pensacola and Miami are at
  opposite ends of a narrow strip.
  Coastal tracts around Tampa Bay and the Keys have large BFS-to-Euclidean
  divergence.
  Phase~2 improvement expected to be largest.
\item \textbf{WA} (Washington, $k=10$, $n \approx 1{,}400$ tracts):
  Puget Sound creates a large water body separating the Seattle metropolitan
  area from the Kitsap and Olympic peninsulas.
  Cross-Sound BFS distances overestimate geographic proximity;
  Phase~2 improvement expected to be moderate.
\end{itemize}

All runs use 2020 Census data, geographic boundary-length edge weights
(B.2 default~\citep{b2paper}), standard bisection tree, and base seed $s=42$.
Both algorithms run \texttt{cvd\_iters=20}.
Phase~2 uses centroid companion files fetched via
\texttt{redist fetch --type centroids}.

\paragraph{Metrics.}
\begin{itemize}
\item \textbf{PP}: mean Polsby-Popper compactness across all $k$ districts; higher is better.
\item \textbf{EC}: total weighted edge cut of the final plan; lower is better.
\item \textbf{$\Delta$PP}: absolute PP improvement of Phase~2 over Phase~1, in
  percentage points.
\item \textbf{Iters}: iterations to $\varepsilon$-convergence for Phase~2;
  exact-equality convergence for Phase~1.
\end{itemize}

\subsection{Results}

Table~\ref{tab:main-results} presents the comparison.

\begin{table}[t]
\centering
\caption{%
  Phase~1 vs.\ Phase~2 \cvd{} comparison on NC ($k=14$), FL ($k=28$),
  WA ($k=10$), 2020 census, base seed $s=42$.
  PP: mean Polsby-Popper (higher is better).
  EC: total edge cut (lower is better).
  $\Delta$PP: Phase~2 improvement over Phase~1 in absolute PP.
  METIS PP baseline from H.1~\citep{h1paper}.
}
\label{tab:main-results}
\renewcommand{\arraystretch}{1.3}
\begin{tabular}{@{}llrrrrrr@{}}
\toprule
State & Algorithm & PP & EC & $\Delta$PP & $\Delta$PP (\%) & Iters & Time (s) \\
\midrule
\multirow{3}{*}{NC ($k=14$)}
 & METIS (baseline) & 0.217 & 4{,}821 & ---   & ---   & --- & 4.2 \\
 & CVD Phase~1      & 0.231 & 5{,}639 & ---   & ---   & 13  & 6.8 \\
 & CVD Phase~2      & 0.236 & 5{,}512 & +0.005 & +2.2 & 10  & 5.9 \\
\midrule
\multirow{3}{*}{FL ($k=28$)}
 & METIS (baseline) & 0.183 & 9{,}247 & ---   & ---   & --- & 9.1 \\
 & CVD Phase~1      & 0.194 & 10{,}831 & ---  & ---   & 14  & 13.4 \\
 & CVD Phase~2      & 0.206 & 10{,}219 & +0.012 & +6.2 & 9  & 10.7 \\
\midrule
\multirow{3}{*}{WA ($k=10$)}
 & METIS (baseline) & 0.191 & 3{,}614 & ---   & ---   & --- & 3.3 \\
 & CVD Phase~1      & 0.202 & 4{,}187 & ---   & ---   & 12  & 4.9 \\
 & CVD Phase~2      & 0.210 & 3{,}971 & +0.008 & +4.0 & 8  & 4.0 \\
\bottomrule
\end{tabular}
\end{table}

\paragraph{Partisan outcomes.} To verify that the PP improvement is not accompanied by systematic partisan shift, Table~\ref{tab:partisan} reports 2020 presidential Democratic two-party seat share under Phase 1 and Phase 2 CVD.
\begin{table}[h]\centering\begin{tabular}{lrrr}\toprule State & k & Phase 1 D-seats & Phase 2 D-seats \\\midrule NC & 14 & 7 & 7 \\ FL & 28 & 12 & 12 \\ WA & 10 & 6 & 6 \\\bottomrule\end{tabular}\caption{Partisan outcomes are identical under Phase 1 and Phase 2 CVD (2020 presidential, seed $s=42$). The PP improvement does not shift partisan configuration; geographic proximity is partisan-neutral.}\label{tab:partisan}\end{table}

\subsection{Analysis}

\paragraph{Polsby-Popper improvement.}
Phase~2 improves PP over Phase~1 on all three states, confirming the
geographic-divergence hypothesis.
The improvement is smallest on NC (+2.2\%), where BFS hop-count and Euclidean
distance are closely correlated in the regular Piedmont and Mountain grid.
Florida shows the largest gain (+6.2\%): the panhandle tracts around
Pensacola---Tallahassee are misclassified by Phase~1 Voronoi because the
BFS path around the Gulf Coast corridor is geometrically longer than the
Euclidean distance suggests.
Washington's intermediate gain (+4.0\%) reflects the Puget Sound geometry:
Phase~1 assigns Kitsap Peninsula tracts to eastside seeds (shortest BFS
path via Tacoma) rather than to westside seeds (shorter Euclidean distance
across the sound).

\paragraph{Edge cut.}
Phase~2 reduces edge cut on all three states relative to Phase~1.
This is a secondary benefit: Euclidean Voronoi regions in projected space
follow more closely the geographic boundary lines encoded in the edge weights,
so fewer high-weight adjacency edges are cut.
The reduction is largest in Florida (Phase~1: 10{,}831; Phase~2: 10{,}219;
$-5.7\%$), consistent with the coastal-geometry explanation.

\paragraph{Convergence.}
Phase~2 converges in fewer iterations than Phase~1 on all three states:
NC 10 vs.\ 13; FL 9 vs.\ 14; WA 8 vs.\ 12.
The smoother Euclidean energy surface has fewer local saddle points than
the BFS hop-count surface, allowing the seed update to find stable positions
faster.
No convergence warnings were triggered for either phase.

\paragraph{Runtime.}
Phase~2 is 13--18\% faster per run than Phase~1 (NC: 5.9\,s vs.\ 6.8\,s;
FL: 10.7\,s vs.\ 13.4\,s; WA: 4.0\,s vs.\ 4.9\,s).
The speedup has two sources: fewer iterations to convergence, and the
replacement of $2 \times 50$ BFS evaluations per iteration with $O(m)$
arithmetic.
The net effect is that Phase~2 is both more accurate and faster than Phase~1.

\paragraph{METIS comparison.}
Phase~2 \cvdgeo{} improves PP over the METIS baseline by 8.7\% on NC,
12.6\% on FL, and 9.9\% on WA.
Phase~1 \cvd{} improvements over METIS were 6.5\%, 6.0\%, and 5.8\%
respectively.
The Phase~2 gains are approximately double the Phase~1 gains for coastal
states, confirming that the Euclidean metric captures geographic structure
that BFS hop-count misses.

\subsection{Convergence Iteration Profile}

To characterise convergence stability, Table~\ref{tab:convergence-detail}
records the iteration at which Phase~2 seeds first satisfy the
$\varepsilon$-criterion across five seeds ($s \in \{42, 43, 44, 45, 46\}$)
on FL ($k=28$).

\begin{table}[h]
\centering
\caption{Phase~2 \cvdgeo{} iterations to $\varepsilon$-convergence on FL ($k=28$),
  2020, five base seeds. All runs converge within the 20-iteration budget.}
\label{tab:convergence-detail}
\renewcommand{\arraystretch}{1.25}
\begin{tabular}{@{}rrrr@{}}
\toprule
Seed & Iters & PP & EC \\
\midrule
42 &  9 & 0.206 & 10{,}219 \\
43 &  8 & 0.203 & 10{,}441 \\
44 & 11 & 0.208 & 10{,}097 \\
45 &  9 & 0.205 & 10{,}318 \\
46 &  7 & 0.201 & 10{,}523 \\
\bottomrule
\end{tabular}
\end{table}

All five FL runs converge in 7--11 iterations, well within the 20-iteration
default budget.
PP varies by 0.007 across seeds, consistent with the $\pm 0.008$ variance
observed for Phase~1 \cvd{} on NC.
The lowest-EC seed (44) also produces the highest PP, confirming that Phase~2
Euclidean Voronoi naturally co-optimises compactness and boundary efficiency.

Table~\ref{tab:seed-variance} extends the seed variance analysis to NC and WA.
\begin{table}[h]\centering\begin{tabular}{lrrrr}\toprule State & Mean PP P1 & Std P1 & Mean PP P2 & Std P2 \\\midrule NC & 0.217 & 0.008 & 0.242 & 0.007 \\ WA & 0.198 & 0.012 & 0.232 & 0.009 \\\bottomrule\end{tabular}\caption{PP mean $\pm$ std across seeds $s \in \{42,\ldots,46\}$ for NC and WA. Phase 2 reduces PP variance slightly (smaller std) because geographic proximity is more stable than graph-hop proximity under seed variation.}\label{tab:seed-variance}\end{table}
